\IfFileExists{stacks-project.cls}{%
\documentclass{stacks-project}
}{%
\documentclass{amsart}
}

% For dealing with references we use the comment environment
\usepackage{verbatim}
\newenvironment{reference}{\comment}{\endcomment}
%\newenvironment{reference}{}{}
\newenvironment{slogan}{\comment}{\endcomment}
\newenvironment{history}{\comment}{\endcomment}

% For commutative diagrams we use Xy-pic
\usepackage[all]{xy}

% We use 2cell for 2-commutative diagrams.
\xyoption{2cell}
\UseAllTwocells

% We use multicol for the list of chapters between chapters
\usepackage{multicol}

% This is generally recommended for better output
\usepackage{lmodern}
\usepackage[T1]{fontenc}

% For cross-file-references
\usepackage{xr-hyper}

% Package for hypertext links:
\usepackage{hyperref}

% For any local file, say "hello.tex" you want to link to please
% use \externaldocument[hello-]{hello}
\externaldocument[introduction-]{introduction}
\externaldocument[conventions-]{conventions}
\externaldocument[sets-]{sets}
\externaldocument[categories-]{categories}
\externaldocument[topology-]{topology}
\externaldocument[sheaves-]{sheaves}
\externaldocument[sites-]{sites}
\externaldocument[stacks-]{stacks}
\externaldocument[fields-]{fields}
\externaldocument[algebra-]{algebra}
\externaldocument[brauer-]{brauer}
\externaldocument[homology-]{homology}
\externaldocument[derived-]{derived}
\externaldocument[simplicial-]{simplicial}
\externaldocument[more-algebra-]{more-algebra}
\externaldocument[smoothing-]{smoothing}
\externaldocument[modules-]{modules}
\externaldocument[sites-modules-]{sites-modules}
\externaldocument[injectives-]{injectives}
\externaldocument[cohomology-]{cohomology}
\externaldocument[sites-cohomology-]{sites-cohomology}
\externaldocument[dga-]{dga}
\externaldocument[dpa-]{dpa}
\externaldocument[sdga-]{sdga}
\externaldocument[hypercovering-]{hypercovering}
\externaldocument[schemes-]{schemes}
\externaldocument[constructions-]{constructions}
\externaldocument[properties-]{properties}
\externaldocument[morphisms-]{morphisms}
\externaldocument[coherent-]{coherent}
\externaldocument[divisors-]{divisors}
\externaldocument[limits-]{limits}
\externaldocument[varieties-]{varieties}
\externaldocument[topologies-]{topologies}
\externaldocument[descent-]{descent}
\externaldocument[perfect-]{perfect}
\externaldocument[more-morphisms-]{more-morphisms}
\externaldocument[flat-]{flat}
\externaldocument[groupoids-]{groupoids}
\externaldocument[more-groupoids-]{more-groupoids}
\externaldocument[etale-]{etale}
\externaldocument[chow-]{chow}
\externaldocument[intersection-]{intersection}
\externaldocument[pic-]{pic}
\externaldocument[weil-]{weil}
\externaldocument[adequate-]{adequate}
\externaldocument[dualizing-]{dualizing}
\externaldocument[duality-]{duality}
\externaldocument[discriminant-]{discriminant}
\externaldocument[derham-]{derham}
\externaldocument[local-cohomology-]{local-cohomology}
\externaldocument[algebraization-]{algebraization}
\externaldocument[curves-]{curves}
\externaldocument[resolve-]{resolve}
\externaldocument[models-]{models}
\externaldocument[functors-]{functors}
\externaldocument[equiv-]{equiv}
\externaldocument[pione-]{pione}
\externaldocument[etale-cohomology-]{etale-cohomology}
\externaldocument[proetale-]{proetale}
\externaldocument[relative-cycles-]{relative-cycles}
\externaldocument[more-etale-]{more-etale}
\externaldocument[trace-]{trace}
\externaldocument[crystalline-]{crystalline}
\externaldocument[spaces-]{spaces}
\externaldocument[spaces-properties-]{spaces-properties}
\externaldocument[spaces-morphisms-]{spaces-morphisms}
\externaldocument[decent-spaces-]{decent-spaces}
\externaldocument[spaces-cohomology-]{spaces-cohomology}
\externaldocument[spaces-limits-]{spaces-limits}
\externaldocument[spaces-divisors-]{spaces-divisors}
\externaldocument[spaces-over-fields-]{spaces-over-fields}
\externaldocument[spaces-topologies-]{spaces-topologies}
\externaldocument[spaces-descent-]{spaces-descent}
\externaldocument[spaces-perfect-]{spaces-perfect}
\externaldocument[spaces-more-morphisms-]{spaces-more-morphisms}
\externaldocument[spaces-flat-]{spaces-flat}
\externaldocument[spaces-groupoids-]{spaces-groupoids}
\externaldocument[spaces-more-groupoids-]{spaces-more-groupoids}
\externaldocument[bootstrap-]{bootstrap}
\externaldocument[spaces-pushouts-]{spaces-pushouts}
\externaldocument[spaces-chow-]{spaces-chow}
\externaldocument[groupoids-quotients-]{groupoids-quotients}
\externaldocument[spaces-more-cohomology-]{spaces-more-cohomology}
\externaldocument[spaces-simplicial-]{spaces-simplicial}
\externaldocument[spaces-duality-]{spaces-duality}
\externaldocument[formal-spaces-]{formal-spaces}
\externaldocument[restricted-]{restricted}
\externaldocument[spaces-resolve-]{spaces-resolve}
\externaldocument[formal-defos-]{formal-defos}
\externaldocument[defos-]{defos}
\externaldocument[cotangent-]{cotangent}
\externaldocument[examples-defos-]{examples-defos}
\externaldocument[algebraic-]{algebraic}
\externaldocument[examples-stacks-]{examples-stacks}
\externaldocument[stacks-sheaves-]{stacks-sheaves}
\externaldocument[criteria-]{criteria}
\externaldocument[artin-]{artin}
\externaldocument[quot-]{quot}
\externaldocument[stacks-properties-]{stacks-properties}
\externaldocument[stacks-morphisms-]{stacks-morphisms}
\externaldocument[stacks-limits-]{stacks-limits}
\externaldocument[stacks-cohomology-]{stacks-cohomology}
\externaldocument[stacks-perfect-]{stacks-perfect}
\externaldocument[stacks-introduction-]{stacks-introduction}
\externaldocument[stacks-more-morphisms-]{stacks-more-morphisms}
\externaldocument[stacks-geometry-]{stacks-geometry}
\externaldocument[moduli-]{moduli}
\externaldocument[moduli-curves-]{moduli-curves}
\externaldocument[examples-]{examples}
\externaldocument[exercises-]{exercises}
\externaldocument[guide-]{guide}
\externaldocument[desirables-]{desirables}
\externaldocument[coding-]{coding}
\externaldocument[obsolete-]{obsolete}
\externaldocument[fdl-]{fdl}
\externaldocument[index-]{index}

% Theorem environments.
%
\theoremstyle{plain}
\newtheorem{theorem}[subsection]{Theorem}
\newtheorem{proposition}[subsection]{Proposition}
\newtheorem{lemma}[subsection]{Lemma}

\theoremstyle{definition}
\newtheorem{definition}[subsection]{Definition}
\newtheorem{example}[subsection]{Example}
\newtheorem{exercise}[subsection]{Exercise}
\newtheorem{situation}[subsection]{Situation}

\theoremstyle{remark}
\newtheorem{remark}[subsection]{Remark}
\newtheorem{remarks}[subsection]{Remarks}

\numberwithin{equation}{subsection}

% Macros
%
\def\lim{\mathop{\mathrm{lim}}\nolimits}
\def\colim{\mathop{\mathrm{colim}}\nolimits}
\def\Spec{\mathop{\mathrm{Spec}}}
\def\Hom{\mathop{\mathrm{Hom}}\nolimits}
\def\Ext{\mathop{\mathrm{Ext}}\nolimits}
\def\SheafHom{\mathop{\mathcal{H}\!\mathit{om}}\nolimits}
\def\SheafExt{\mathop{\mathcal{E}\!\mathit{xt}}\nolimits}
\def\Sch{\mathit{Sch}}
\def\Mor{\mathop{\mathrm{Mor}}\nolimits}
\def\Ob{\mathop{\mathrm{Ob}}\nolimits}
\def\Sh{\mathop{\mathit{Sh}}\nolimits}
\def\NL{\mathop{N\!L}\nolimits}
\def\CH{\mathop{\mathrm{CH}}\nolimits}
\def\proetale{{pro\text{-}\acute{e}tale}}
\def\etale{{\acute{e}tale}}
\def\QCoh{\mathit{QCoh}}
\def\Ker{\mathop{\mathrm{Ker}}}
\def\Im{\mathop{\mathrm{Im}}}
\def\Coker{\mathop{\mathrm{Coker}}}
\def\Coim{\mathop{\mathrm{Coim}}}

% Boxtimes
%
\DeclareMathSymbol{\boxtimes}{\mathbin}{AMSa}{"02}

%
% Macros for moduli stacks/spaces
%
\def\QCohstack{\mathcal{QC}\!\mathit{oh}}
\def\Cohstack{\mathcal{C}\!\mathit{oh}}
\def\Spacesstack{\mathcal{S}\!\mathit{paces}}
\def\Quotfunctor{\mathrm{Quot}}
\def\Hilbfunctor{\mathrm{Hilb}}
\def\Curvesstack{\mathcal{C}\!\mathit{urves}}
\def\Polarizedstack{\mathcal{P}\!\mathit{olarized}}
\def\Complexesstack{\mathcal{C}\!\mathit{omplexes}}
% \Pic is the operator that assigns to X its picard group, usage \Pic(X)
% \Picardstack_{X/B} denotes the Picard stack of X over B
% \Picardfunctor_{X/B} denotes the Picard functor of X over B
\def\Pic{\mathop{\mathrm{Pic}}\nolimits}
\def\Picardstack{\mathcal{P}\!\mathit{ic}}
\def\Picardfunctor{\mathrm{Pic}}
\def\Deformationcategory{\mathcal{D}\!\mathit{ef}}
\usepackage{cleveref}
\crefname{originalsection}{original section}{original sections}
\Crefname{originalsection}{Original section}{Original sections}
\makeatletter
\newlabel{cohomology-section-global-RHom}{{0B6A}{source}{}{sourcepointer-1}{}}
\newlabel{cohomology-section-global-RHom@cref}{{[originalsection][0][]0B6A}{[1][0][]source}{}{}{}}
\makeatother
\crefname{originallemma}{original lemma}{original lemmas}
\Crefname{originallemma}{Original lemma}{Original lemmas}
\makeatletter
\newlabel{cohomology-lemma-dual-perfect-complex}{{08DQ}{source}{}{sourcepointer-2}{}}
\newlabel{cohomology-lemma-dual-perfect-complex@cref}{{[originallemma][0][]08DQ}{[1][0][]source}{}{}{}}
\makeatother
\crefname{originallemma}{original lemma}{original lemmas}
\Crefname{originallemma}{Original lemma}{Original lemmas}
\makeatletter
\newlabel{cohomology-lemma-tensor-pseudo-coherent}{{09J3}{source}{}{sourcepointer-3}{}}
\newlabel{cohomology-lemma-tensor-pseudo-coherent@cref}{{[originallemma][0][]09J3}{[1][0][]source}{}{}{}}
\makeatother
\crefname{originallemma}{original lemma}{original lemmas}
\Crefname{originallemma}{Original lemma}{Original lemmas}
\makeatletter
\newlabel{lemma-flat-proper-pseudo-coherent-direct-image-general}{{0CSD}{source}{}{sourcepointer-4}{}}
\newlabel{lemma-flat-proper-pseudo-coherent-direct-image-general@cref}{{[originallemma][0][]0CSD}{[1][0][]source}{}{}{}}
\makeatother
\crefname{originalequation}{original equation}{original equations}
\Crefname{originalequation}{Original equation}{Original equations}
\makeatletter
\newlabel{equation-generator-Pn}{{0CS8}{source}{}{sourcepointer-5}{}}
\newlabel{equation-generator-Pn@cref}{{[originalequation][0][]0CS8}{[1][0][]source}{}{}{}}
\makeatother
\crefname{originalequation}{original equation}{original equations}
\Crefname{originalequation}{Original equation}{Original equations}
\makeatletter
\newlabel{equation-algebra-for-Pn}{{0CS9}{source}{}{sourcepointer-6}{}}
\newlabel{equation-algebra-for-Pn@cref}{{[originalequation][0][]0CS9}{[1][0][]source}{}{}{}}
\makeatother
\crefname{originallemma}{original lemma}{original lemmas}
\Crefname{originallemma}{Original lemma}{Original lemmas}
\makeatletter
\newlabel{lemma-Pn-module-category}{{0BQU}{source}{}{sourcepointer-7}{}}
\newlabel{lemma-Pn-module-category@cref}{{[originallemma][0][]0BQU}{[1][0][]source}{}{}{}}
\makeatother
\crefname{originallemma}{original lemma}{original lemmas}
\Crefname{originallemma}{Original lemma}{Original lemmas}
\makeatletter
\newlabel{dga-lemma-upgrade-tensor-with-complex-derived}{{0CS6}{source}{}{sourcepointer-8}{}}
\newlabel{dga-lemma-upgrade-tensor-with-complex-derived@cref}{{[originallemma][0][]0CS6}{[1][0][]source}{}{}{}}
\makeatother
\crefname{originallemma}{original lemma}{original lemmas}
\Crefname{originallemma}{Original lemma}{Original lemmas}
\makeatletter
\newlabel{lemma-pseudo-coherent-over-algebra}{{0CSF}{source}{}{sourcepointer-9}{}}
\newlabel{lemma-pseudo-coherent-over-algebra@cref}{{[originallemma][0][]0CSF}{[1][0][]source}{}{}{}}
\makeatother
\crefname{originalsection}{original section}{original sections}
\Crefname{originalsection}{Original section}{Original sections}
\makeatletter
\newlabel{dga-section-P-resolutions}{{09KK}{source}{}{sourcepointer-10}{}}
\newlabel{dga-section-P-resolutions@cref}{{[originalsection][0][]09KK}{[1][0][]source}{}{}{}}
\makeatother
\crefname{originallemma}{original lemma}{original lemmas}
\Crefname{originallemma}{Original lemma}{Original lemmas}
\makeatletter
\newlabel{dga-lemma-tensor-with-complex-derived}{{09LX}{source}{}{sourcepointer-11}{}}
\newlabel{dga-lemma-tensor-with-complex-derived@cref}{{[originallemma][0][]09LX}{[1][0][]source}{}{}{}}
\makeatother
\title{Pseudo-coherence on projective space is detected by finitely many twists}
\author{The Stacks Project Authors}
\date{stacks-project:perfect.tex:local-snapshot-20261012; GNU FDL 1.2 or later}
\begin{document}
\maketitle
\setcounter{subsection}{1}
\noindent\textit{Editorial note.} All delivered mathematical text is the author's own, copied from the pinned Stacks Project chapter file perfect.tex; no statement or proof step was rewritten, reordered or removed. The counted span is the author environment and its own proof environment. References inside the span are handed to the assembler as navigation stubs that carry the canonical tag of the cited tag only; whether each of them is load-bearing is recorded on the stub itself.\par\medskip
\noindent\textit{Allowed standard prerequisites (editorial).}\begin{itemize}
\item Every perfect complex is pseudo-coherent and the tensor product of two pseudo-coherent complexes is pseudo-coherent. used at proof lines 8797-8806. Gives the implication from the first condition to the second and the trivial implication from the second to the third.
\item The direct image of a pseudo-coherent complex along a flat proper morphism of finite presentation is pseudo-coherent. invoked at proof lines 8804-8806. Produces the second condition from the first.
\item The derived category of quasi-coherent modules on projective space is equivalent to the derived category of modules over the graded algebra of the projective space, with the generator corresponding to the structure sheaf. invoked at proof lines 8810-8820. Lets the problem be translated into a statement about modules over a ring.
\item The derived Hom over the ambient algebra computes the graded pieces of the tensor product with the generator, and pseudo-coherence descends along the equivalence. invoked at proof lines 8821-8862. Concludes that the given complex is pseudo-coherent from the condition on the twists.
\end{itemize}\medskip
\ifdefined\hypertarget\hypertarget{sourcepointer-1}{}\fi
\noindent\textit{Source reference.} Original section 0B6A: sources/stacks-chow/tags.txt, source lines 1--1. Cited at proof lines 8794-8796 to identify the two expressions and thus the equivalence of conditions (3) and (4); the deduction is inside the counted span.\par\medskip
\ifdefined\hypertarget\hypertarget{sourcepointer-2}{}\fi
\noindent\textit{Source reference.} Original lemma 08DQ: sources/stacks-chow/tags.txt, source lines 1--1. Used at proof lines 8798-8799 to finish the equivalence of conditions (3) and (4); everything else used for it is inside the counted span.\par\medskip
\ifdefined\hypertarget\hypertarget{sourcepointer-3}{}\fi
\noindent\textit{Source reference.} Original lemma 09J3: sources/stacks-chow/tags.txt, source lines 1--1. Invoked at proof lines 8801-8803 to see that the tensor product with a pseudo-coherent object is pseudo-coherent; the implication from the first to the second condition is inside the counted span.\par\medskip
\ifdefined\hypertarget\hypertarget{sourcepointer-4}{}\fi
\noindent\textit{Source reference.} Original lemma 0CSD: sources/stacks-perfect/pinned-original-source.tex, source lines 7764--7764. Used at proof lines 8804-8806 to conclude that the direct image of the tensor product is pseudo-coherent, which is the second condition; the deduction is inside the counted span.\par\medskip
\ifdefined\hypertarget\hypertarget{sourcepointer-5}{}\fi
\noindent\textit{Source reference.} Original equation 0CS8: sources/stacks-perfect/pinned-original-source.tex, source lines 4920--4920. Cited at proof lines 8809-8811 to fix the object P used in the translation to modules; the translation argument is inside the counted span.\par\medskip
\ifdefined\hypertarget\hypertarget{sourcepointer-6}{}\fi
\noindent\textit{Source reference.} Original equation 0CS9: sources/stacks-perfect/pinned-original-source.tex, source lines 4927--4927. Cited with the previous label at proof lines 8809-8811 for the algebra R; the use made of it is inside the counted span.\par\medskip
\ifdefined\hypertarget\hypertarget{sourcepointer-7}{}\fi
\noindent\textit{Source reference.} Original lemma 0BQU: sources/stacks-perfect/pinned-original-source.tex, source lines 4961--4961. Invoked at proof lines 8812-8816 to obtain the equivalence and the object M mapping to K; the rest of the argument is inside the counted span.\par\medskip
\ifdefined\hypertarget\hypertarget{sourcepointer-8}{}\fi
\noindent\textit{Source reference.} Original lemma 0CS6: sources/stacks-chow/tags.txt, source lines 1--1. Cited at proof lines 8818-8822 for the displayed isomorphism; the use of it to reach the fifth condition is inside the counted span.\par\medskip
\ifdefined\hypertarget\hypertarget{sourcepointer-9}{}\fi
\noindent\textit{Source reference.} Original lemma 0CSF: sources/stacks-perfect/pinned-original-source.tex, source lines 8687--8687. Invoked at proof lines 8855-8860 to conclude that the module M is pseudo-coherent, hence so is K; the identification of the derived Hom that makes it applicable is inside the counted span.\par\medskip
\ifdefined\hypertarget\hypertarget{sourcepointer-10}{}\fi
\noindent\textit{Source reference.} Original section 09KK: sources/stacks-chow/tags.txt, source lines 1--1. A navigation pointer in the algebra side of the argument; the descent itself is inside the counted span.\par\medskip
\ifdefined\hypertarget\hypertarget{sourcepointer-11}{}\fi
\noindent\textit{Source reference.} Original lemma 09LX: sources/stacks-chow/tags.txt, source lines 1--1. A navigation pointer used alongside the previous one; the use made of it is inside the counted span.\par\medskip
\begin{lemma}
\label{lemma-pseudo-coherent-on-projective-space}
Let $A$ be a ring. Let $n \geq 0$. Let
$K \in D_\QCoh(\mathcal{O}_{\mathbf{P}^n_A})$.
The following are equivalent
\begin{enumerate}
\item $K$ is pseudo-coherent,
\item $R\Gamma(\mathbf{P}^n_A, E \otimes^\mathbf{L} K)$ is a pseudo-coherent
object of $D(A)$ for each pseudo-coherent object $E$ of
$D(\mathcal{O}_{\mathbf{P}^n_A})$,
\item $R\Gamma(\mathbf{P}^n_A, E \otimes^\mathbf{L} K)$ is a pseudo-coherent
object of $D(A)$ for each perfect object $E$ of
$D(\mathcal{O}_{\mathbf{P}^n_A})$,
\item $R\Hom_{\mathbf{P}^n_A}(E, K)$ is a pseudo-coherent
object of $D(A)$ for each perfect object $E$ of
$D(\mathcal{O}_{\mathbf{P}^n_A})$,
\item $R\Gamma(\mathbf{P}^n_A,
K \otimes^\mathbf{L} \mathcal{O}_{\mathbf{P}^n_A}(d))$ is pseudo-coherent
object of $D(A)$ for $d = 0, 1, \ldots, n$.
\end{enumerate}
\end{lemma}

\begin{proof}
Recall that
$$
R\Hom_{\mathbf{P}^n_A}(E, K) =
R\Gamma(\mathbf{P}^n_A, R\SheafHom_{\mathcal{O}_{\mathbf{P}^n_A}}(E, K))
$$
by definition, see Cohomology, Section \ref{cohomology-section-global-RHom}.
Thus parts (4) and (3) are equivalent by
Cohomology, Lemma \ref{cohomology-lemma-dual-perfect-complex}.

\medskip\noindent
Since every perfect complex is pseudo-coherent, it is clear that
(2) implies (3).

\medskip\noindent
Assume (1) holds. Then $E \otimes^\mathbf{L} K$ is pseudo-coherent
for every pseudo-coherent $E$, see
Cohomology, Lemma \ref{cohomology-lemma-tensor-pseudo-coherent}.
By Lemma \ref{lemma-flat-proper-pseudo-coherent-direct-image-general}
the direct image of such a pseudo-coherent complex is pseudo-coherent
and we see that (2) is true.

\medskip\noindent
Part (3) implies (5) because we can take $E = \mathcal{O}_{\mathbf{P}^n_A}(d)$
for $d = 0, 1, \ldots, n$.

\medskip\noindent
To finish the proof we have to show that (5) implies (1).
Let $P$ be as in (\ref{equation-generator-Pn}) and
$R$ as in (\ref{equation-algebra-for-Pn}).
By Lemma \ref{lemma-Pn-module-category} we have an equivalence
$$
- \otimes^\mathbf{L}_R P :
D(R) \longrightarrow D_\QCoh(\mathcal{O}_{\mathbf{P}^n_A})
$$
Let $M \in D(R)$ be an object such that $M \otimes^\mathbf{L} P = K$.
By Differential Graded Algebra, Lemma
\ref{dga-lemma-upgrade-tensor-with-complex-derived}
there is an isomorphism
$$
R\Hom(R, M) = R\Hom_{\mathbf{P}^n_A}(P, K)
$$
in $D(A)$. Arguing as above we obtain
$$
R\Hom_{\mathbf{P}^n_A}(P, K) = R\Gamma(\mathbf{P}^n_A,
R\SheafHom_{\mathcal{O}_{\mathbf{P}^n_A}}(E, K)) =
R\Gamma(\mathbf{P}^n_A,
P^\vee \otimes^\mathbf{L}_{\mathcal{O}_{\mathbf{P}^n_A}} K).
$$
Using that $P^\vee$ is the direct sum of
$\mathcal{O}_{\mathbf{P}^n_A}(d)$ for $d = 0, 1, \ldots, n$
and (5) we conclude $R\Hom(R, M)$ is pseudo-coherent as a complex
of $A$-modules. Of course $M = R\Hom(R, M)$ in $D(A)$.
Thus $M$ is pseudo-coherent as a complex of $A$-modules.
By Lemma \ref{lemma-pseudo-coherent-over-algebra}
we may represent $M$ by a bounded above complex $F^\bullet$
of finite free $R$-modules. Then
$F^\bullet = \bigcup_{p \geq 0} \sigma_{\geq p}F^\bullet$
is a filtration which shows that $F^\bullet$ is a differential
graded $R$-module with property (P), see
Differential Graded Algebra, Section \ref{dga-section-P-resolutions}.
Hence $K = M \otimes^\mathbf{L}_R P$ is represented
by $F^\bullet \otimes_R P$ (follows from the construction of the
derived tensor functor, see for example the proof of
Differential Graded Algebra, Lemma \ref{dga-lemma-tensor-with-complex-derived}).
Since $F^\bullet \otimes_R P$
is a bounded above complex whose terms are direct sums of
copies of $P$ we conclude that the lemma is true.
\end{proof}

\end{document}
